% GENERATED FILE. Edit canonical lessons, then run npm run generate:books.
% canonical-source-hash: fnv1a64:af0d6d57636bd1fa
% canonical-lessons: MR-C171-kvacit, MR-C171-mag, MR-C171-sevti, MR-C171-ajun, MR-C171-acanak

\chapter{Rarely, Then, Finally, Still, Suddenly}
\label{ch:mr-rarely-then-finally-still-suddenly}
\hlchaptermodality{171}

\begin{chapteropening}
\textbf{By the end of this chapter:} \emph{I can say rarely, then, finally, still and suddenly.}

Complete the last lesson of chapter 171: I can say rarely, then, finally, still and suddenly.
\end{chapteropening}

\section[\texorpdfstring{\mr{क्वचित} (kvacit)}{क्वचित (kvacit)}]{\mr{क्वचित} (kvacit) — rarely}

\label{lesson:MR-C171-kvacit}



\begin{quote}
Before the new one: say the Marathi for patience, restraint, then the Marathi for luck, fate.
\end{quote}

\subsection*{You'll want to know: \mr{क्वचित}}
\textbf{\mr{क्वचित}} — \emph{kvacit} — \textquotedblleft{}rarely\textquotedblright{}.

\begin{cousinweb}
\textbf{\mr{क्वचित}} (\emph{kvacit}) is Sanskrit \emph{kvacit}, \textquotedblleft{}somewhere, at times\textquotedblright{}, a tatsama: \emph{kva}, \textquotedblleft{}where\textquotedblright{}, with \emph{-cit}, "some-". CDIAL has no entry for it.
\end{cousinweb}

\begin{grammarlens}[title={What you've built}]
One more word for this chapter.
\end{grammarlens}

\subsection*{Your turn}
\begin{itemize}
\raggedright
  \item \emph{Say it:} \emph{kvacit}
  \item \emph{Say it:} \emph{kvacit}, once more
  \item \emph{Say it:} say \emph{kvacit}
  \item \emph{Recall:} say the Marathi for patience, restraint, then the Marathi for luck, fate, then say \emph{kvacit} again
\end{itemize}

\begin{tcolorbox}[breakable,colback=teal!4,colframe=teal!35!black,title={Before you move on}]
What does \mr{क्वचित} mean? (\textquotedblleft{}Rarely\textquotedblright{}.) Say it once more.
\end{tcolorbox}
\section[\texorpdfstring{\mr{मग} (mag)}{मग (mag)}]{\mr{मग} (mag) — then, after that}

\label{lesson:MR-C171-mag}



\begin{quote}
Before the new one: say the Marathi for luck, fate, then the Marathi for rarely.
\end{quote}

\subsection*{You'll want to know: \mr{मग}}
\textbf{\mr{मग}} — \emph{mag} — \textquotedblleft{}then, after that\textquotedblright{}.

\begin{cousinweb}
\textbf{\mr{मग}} (\emph{mag}) is inherited. CDIAL lists Marathi \emph{mag}, \textquotedblleft{}then, after\textquotedblright{}, as an adverbial use of Sanskrit \emph{mārga}, \textquotedblleft{}track, road\textquotedblright{} (CDIAL 10071).
\end{cousinweb}

\begin{grammarlens}[title={What you've built}]
One more word for this chapter.
\end{grammarlens}

\subsection*{Your turn}
\begin{itemize}
\raggedright
  \item \emph{Say it:} \emph{mag}
  \item \emph{Say it:} \emph{mag}, once more
  \item \emph{Say it:} say \emph{mag}
  \item \emph{Recall:} say the Marathi for luck, fate, then the Marathi for rarely, then say \emph{mag} again
\end{itemize}

\begin{tcolorbox}[breakable,colback=teal!4,colframe=teal!35!black,title={Before you move on}]
What does \mr{मग} mean? (\textquotedblleft{}Then, after that\textquotedblright{}.) Say it once more.
\end{tcolorbox}
\section[\texorpdfstring{\mr{शेवटी} (śevṭī)}{शेवटी (śevṭī)}]{\mr{शेवटी} (śevṭī) — finally, in the end}

\label{lesson:MR-C171-sevti}



\begin{quote}
Before the new one: say the Marathi for rarely, then the Marathi for then.
\end{quote}

\subsection*{You'll want to know: \mr{शेवटी}}
\textbf{\mr{शेवटी}} — \emph{śevṭī} — \textquotedblleft{}finally, in the end\textquotedblright{}.

\begin{cousinweb}
The origin of \textbf{\mr{शेवटी}} (\emph{śevṭī}) is not settled. Neither CDIAL nor DEDR, searched by form and by meaning, has an entry for it.
\end{cousinweb}

\begin{grammarlens}[title={What you've built}]
One more word for this chapter.
\end{grammarlens}

\subsection*{Your turn}
\begin{itemize}
\raggedright
  \item \emph{Say it:} \emph{śevṭī}
  \item \emph{Say it:} \emph{śevṭī}, once more
  \item \emph{Say it:} say \emph{śevṭī}
  \item \emph{Recall:} say the Marathi for rarely, then the Marathi for then, then say \emph{śevṭī} again
\end{itemize}

\begin{tcolorbox}[breakable,colback=teal!4,colframe=teal!35!black,title={Before you move on}]
What does \mr{शेवटी} mean? (\textquotedblleft{}Finally, in the end\textquotedblright{}.) Say it once more.
\end{tcolorbox}
\section[\texorpdfstring{\mr{अजून} (ajūn)}{अजून (ajūn)}]{\mr{अजून} (ajūn) — still, yet}

\label{lesson:MR-C171-ajun}



\begin{quote}
Before the new one: say the Marathi for then, then the Marathi for finally.
\end{quote}

\subsection*{You'll want to know: \mr{अजून}}
\textbf{\mr{अजून}} — \emph{ajūn} — \textquotedblleft{}still, yet\textquotedblright{}.

\begin{cousinweb}
\textbf{\mr{अजून}} (\emph{ajūn}) is a tadbhava, a word inherited and reshaped by sound change. It comes from \emph{*adyāhnaḥ}, \textquotedblleft{}this very day\textquotedblright{}, a form CDIAL reconstructs from its descendants (CDIAL 244). It is \emph{adya}, \textquotedblleft{}today\textquotedblright{}, and \emph{ahnaḥ}, \textquotedblleft{}of the day\textquotedblright{}. CDIAL glosses the Marathi word \textquotedblleft{}even now\textquotedblright{}.
\end{cousinweb}

\begin{grammarlens}[title={What you've built}]
One more word for this chapter.
\end{grammarlens}

\subsection*{Your turn}
\begin{itemize}
\raggedright
  \item \emph{Say it:} \emph{ajūn}
  \item \emph{Say it:} \emph{ajūn}, once more
  \item \emph{Say it:} say \emph{ajūn}
  \item \emph{Recall:} say the Marathi for then, then the Marathi for finally, then say \emph{ajūn} again
\end{itemize}

\begin{tcolorbox}[breakable,colback=teal!4,colframe=teal!35!black,title={Before you move on}]
What does \mr{अजून} mean? (\textquotedblleft{}Still, yet\textquotedblright{}.) Say it once more.
\end{tcolorbox}
\section[\texorpdfstring{\mr{अचानक} (acānak)}{अचानक (acānak)}]{\mr{अचानक} (acānak) — suddenly}

\label{lesson:MR-C171-acanak}



\begin{quote}
Before the new one: say the Marathi for finally, then the Marathi for still.
\end{quote}

\subsection*{You'll want to know: \mr{अचानक}}
\textbf{\mr{अचानक}} — \emph{acānak} — \textquotedblleft{}suddenly\textquotedblright{}.

\begin{cousinweb}
The origin of \textbf{\mr{अचानक}} (\emph{acānak}) is uncertain. CDIAL has no entry for it.
\end{cousinweb}

\begin{grammarlens}[title={What you've built}]
One more word for this chapter.
\end{grammarlens}

\subsection*{Your turn}
\begin{itemize}
\raggedright
  \item \emph{Say it:} \emph{acānak}
  \item \emph{Say it:} \emph{acānak}, once more
  \item \emph{Say it:} say \emph{acānak}
  \item \emph{Recall:} say the Marathi for finally, then the Marathi for still, then say \emph{acānak} again
\end{itemize}

\begin{tcolorbox}[breakable,colback=teal!4,colframe=teal!35!black,title={Before you move on}]
What does \mr{अचानक} mean? (\textquotedblleft{}Suddenly\textquotedblright{}.) Say it once more.
\end{tcolorbox}
